\documentclass[../../../include/open-logic-section]{subfiles}

\begin{document}

\olfileid{sth}{spine}{rank}
\olsection{رتبه}

اکنون که مراحل را به‌منزلهٔ~$V_\alpha$ها تعریف کرده‌ایم و می‌دانیم هر
مجموعه زیرمجموعهٔ مرحله‌ای است، می‌توانیم \emph{رتبهٔ} یک مجموعه را تعریف
کنیم. به‌طور شهودی، رتبهٔ~$A$ نخستین لحظه‌ای است که~$A$ در آن تشکیل
می‌شود. به‌بیان دقیق‌تر:

\begin{defn}\ollabel{defnsetrank}
برای هر مجموعهٔ~$A$، $\setrank{A}$ کمترین عدد ترتیبی~$\alpha$ است
به‌طوری‌که~$A\subseteq V_\alpha$.
\end{defn}
\begin{prop}\ollabel{ranksexist}
	برای هر~$A$، $\setrank{A}$ وجود دارد.
\end{prop}
\begin{proof}
	به‌عنوان تمرین واگذار می‌شود.
\end{proof}
\begin{prob}
	\olref[sth][spine][rank]{ranksexist} را اثبات کنید.
\end{prob}
\noindent
خوش‌ترتیبی رتبه‌ها به ما اجازه می‌دهد چند نتیجهٔ مهم را اثبات کنیم:
\begin{prop}\ollabel{valphalowerrank}
برای هر عدد ترتیبی~$\alpha$،
$V_\alpha=\Setabs{x}{\setrank{x}\in\alpha}$.
\end{prop}

\begin{proof}
اگر~$\setrank{x}\in\alpha$، آنگاه
$x\subseteq V_{\setrank{x}}\in V_\alpha$؛ پس چون~$V_\alpha$ توانمند
است، با چند بار استفاده از~\olref[Valphabasic]{Valphabasicprops} نتیجه
می‌شود~$x\in V_\alpha$. برعکس، اگر~$x\in V_\alpha$، آنگاه
$x\subseteq V_\alpha$ و بنابراین~$\setrank{x}\leq\alpha$؛ و استقرای
ترامتناهی ساده‌ای نشان می‌دهد~$x\notin V_{\setrank{x}}$. پس
$\setrank{x}\neq\alpha$ و در نتیجه~$\setrank{x}\in\alpha$.
\end{proof}
\begin{prob}
	استقرای ترامتناهی ساده‌ای را که در
	\olref[sth][spine][rank]{valphalowerrank} ذکر شد کامل کنید.
\end{prob}

\begin{prop}\ollabel{rankmemberslower}
اگر~$B\in A$، آنگاه~$\setrank{B}\in\setrank{A}$.
\end{prop}

\begin{proof}
$A\subseteq V_{\setrank{A}}=
\Setabs{x}{\setrank{x}\in\setrank{A}}$ بنا بر~\olref{valphalowerrank}.
\end{proof}
\noindent
با استفاده از این واقعیت می‌توانیم نتیجه‌ای به دست آوریم که نوعی استقرا
روی \emph{همهٔ مجموعه‌ها} را ممکن می‌سازد:

\begin{thm}[طرحوارهٔ استقرای $\in$]
برای هر فرمول~$\phi$:
\[
	\forall A((\forall x \in A)\phi(x) \lif \phi(A)) \lif \forall A \phi(A).
\]
\end{thm}

\begin{proof}
عکس نقیض را اثبات می‌کنیم. فرض کنید~$\lnot\forall A\phi(A)$. بنا بر
استقرای ترامتناهی
(\olref[ordinals][basic]{ordinductionschema})، یک مجموعهٔ فاقد~$\phi$ با
کمترین رتبهٔ ممکن وجود دارد؛ یعنی~$A$ای هست به‌طوری‌که
$\lnot\phi(A)$ و
$\forall x(\setrank{x}\in\setrank{A}\lif\phi(x))$. حال اگر~$x\in A$،
آنگاه بنا بر~\olref{rankmemberslower}،
$\setrank{x}\in\setrank{A}$، و ازاین‌رو~$\phi(x)$؛ یعنی
$(\forall x\in A)\phi(x)\land\lnot\phi(A)$.
\end{proof}\noindent
شرح غیرصوریِ این نتیجهٔ نیرومند چنین است. می‌گوییم~$\phi$
\emph{وراثتی} است اگر و تنها اگر هرگاه هر~!!{element} یک مجموعه دارای
$\phi$ باشد، خود مجموعه نیز دارای~$\phi$ باشد. استقرای~$\in$ می‌گوید:
اگر~$\phi$ وراثتی باشد، هر مجموعه دارای~$\phi$ است.

برای پایان‌دادن به بحث رتبه‌ها (دست‌کم فعلاً)، چند ادعایی را اثبات می‌کنیم
که پیش‌تر چند بار نویدشان را داده بودیم.

\begin{prop}\ollabel{ranksupstrict}
$\setrank{A}=\supstrict_{x\in A}\setrank{x}$.
\end{prop}

\begin{proof}
بگذارید~$\alpha=\supstrict_{x\in A}\setrank{x}$. بنا بر
\olref{rankmemberslower}، $\alpha\leq\setrank{A}$. اما اگر~$x\in A$،
آنگاه~$\setrank{x}\in\alpha$؛ پس بنا بر~\olref{valphalowerrank}،
$x\in V_\alpha$ و بنابراین~$A\subseteq V_\alpha$؛ یعنی
$\setrank{A}\leq\alpha$. پس~$\setrank{A}=\alpha$.
\end{proof}

\begin{cor}\ollabel{ordsetrankalpha}
برای هر عدد ترتیبی~$\alpha$، $\setrank{\alpha}=\alpha$.
\end{cor}

\begin{proof}
برای استقرای ترامتناهی فرض کنید برای همهٔ~$\beta\in\alpha$،
$\setrank{\beta}=\beta$. اکنون
$\setrank{\alpha}=\supstrict_{\beta\in\alpha}\setrank{\beta}=
\supstrict_{\beta\in\alpha}\beta=\alpha$ بنا بر
\olref{ranksupstrict}.
%	نخست توجه کنید $\setrank{\alpha}\neq\beta$ برای هر $\beta\in\alpha$.
%	زیرا در غیر این صورت بنا بر \olref{rankmemberslower} می‌داشتیم
%	$\beta=\setrank{\beta}\in\setrank{\alpha}=\beta$، یعنی
%	$\beta\in\beta$، که تناقض است.
%
%	اکنون توجه کنید $\alpha\subseteq V_\alpha$. زیرا
%	$\alpha=\Setabs{\beta}{\beta\in\alpha}$ بنا بر
%	\olref[sfr][ordinals][basic]{ordissetofsmallerord}. و اگر
%	$\beta\in\alpha$، آنگاه
%	$\beta\subseteq V_{\setrank{\beta}}=V_\beta\in V_\alpha$، و ازاین‌رو
%	با دو بار استفاده از \olref[sfr][spine][Valphabasic]{Valphabasicprops}،
%	$\beta\in V_\alpha$.
\end{proof}

سرانجام، در اینجا اثباتی کوتاه از نتیجه‌ای می‌آوریم که در پایان
\olref[foundation]{sec} وعده داده شد: اینکه~$\ZFminus$ شرطیِ
$\emph{قاعده‌مندی}\Rightarrow\emph{بنیاد}$ را اثبات می‌کند. (توجه کنید
مفهوم «رتبه» و~\olref{rankmemberslower} در این اثبات در دسترس‌اند، زیرا
همان‌طور که در آغاز این بخش گفتیم، می‌توان آن‌ها را با استفاده از
$\ZFminus+\text{قاعده‌مندی}$ ارائه کرد.)

\begin{prop}[در $\ZFminus + \text{قاعده‌مندی}$]\ollabel{zfminusregularityfoundation}
بنیاد برقرار است.
\end{prop}

\begin{proof}
$A\neq\emptyset$ را ثابت بگیرید و~$B\in A$ را با کمترین رتبهٔ ممکن
انتخاب کنید. اگر~$c\in B$، آنگاه بنا بر~\olref{rankmemberslower}،
$\setrank{c}\in\setrank{B}$؛ پس به‌سبب انتخاب~$B$، $c\notin A$.
\end{proof}

\end{document}
